\renewcommand{\theposttitle}{I.D.E.A. Newsletter 9: Storytelling and Comfort Zones}
\renewcommand{\thepostdate}{2022-07-22}
\vspace*{\fill}
\begin{center}
\includegraphics[width=\textwidth, height=0.55\textheight, keepaspectratio]{C:/Users/James/Documents/GitHub/VisualIdeasBook/data/images/4d0110068e4ee382f4b22ff2562e5d05.jpeg}
\end{center}
\vspace{2mm}
\begin{center}
\begin{minipage}{0.92\textwidth}
\small
The highlight of my week was when my friend @EvansNifty posted this visual demonstrating my philosophy on storytelling.

If you want your content to stand out, involve the audience as active participants of your work.Telling a good story is as if you're drawing a circle for your audience.

But instead of forming the circle completely by yourself, you want to leave it slightly incomplete.

That small gap at the end is for the audience to complete themselves.

Don't give the audience the answer (4). Engage them with the question (2+2).
\end{minipage}
\end{center}
\vspace*{\fill}